\clearpage\thispagestyle{empty}
\pdfbookmark[0]{Quellen zu Gauß an Bode (17. März und 15. August 1817)}{bodemaraug1817-sources}
\begin{center}\textbf{Quellen zu Gauß an Bode (17. März und 15. August 1817)}\end{center}
\noindent Diese KI-Übertragung ins heutige Deutsch folgt den Briefen von Carl Friedrich Gauß an Johann Elert Bode in den \href{https://gdz.sub.uni-goettingen.de/id/PPN236007467}{\textit{Werken}, Band~VI} (Göttingen, 1874), S.~399--402. Der Erstabdruck erschien unter dem gemeinsamen Titel \textit{Astronomische Beobachtungen und Bemerkungen, vom Hrn. Prof. Ritter Gauß in Göttingen, aus Briefen desselben} im \textit{Astronomischen Jahrbuch für 1820} (Berlin, 1817), S.~\href{https://api.digitale-sammlungen.de/iiif/image/v2/bsb10538349_00205/full/1600,/0/default.jpg}{201}--\href{https://api.digitale-sammlungen.de/iiif/image/v2/bsb10538349_00209/full/1600,/0/default.jpg}{205}. Die Links zum Erstabdruck führen zum Exemplar der Bayerischen Staatsbibliothek.
\par\smallskip\noindent Auf der Anfangsseite des Jahrbuchs stehen zwei Anmerkungen zum Abdruck. Die erste folgt auf die Erwähnung der Ephemeride und lautet „Sie folgt weiter unten.“ Die zweite folgt auf Gauß’ Frage nach dem winterlichen Himmel und lautet „Siehe meine Beobachtungen auf S.~90 und den folgenden Seiten.“ Beide Anmerkungen fehlen in der Werkausgabe. Die im Brief erwähnte Ephemeride gehört nicht zum dort wiedergegebenen Text.
\par\smallskip\noindent Auf S.~\href{https://api.digitale-sammlungen.de/iiif/image/v2/bsb10538349_00206/full/1600,/0/default.jpg}{202} verwendet das Jahrbuch kurze Striche für die Eintritte bei der Mondfinsternis und schreibt den östlichen Versatz als $1''{,}90$ in Zeit. Die Werkausgabe hat Punktreihen und $1^{\mathrm s}90$. Die Einträge 16,16 für den 11.~April und 49,60 für den 7.~Mai auf S.~\href{https://api.digitale-sammlungen.de/iiif/image/v2/bsb10538349_00207/full/1600,/0/default.jpg}{203} und der aus den Sonnenbeobachtungen abgeleitete Wert $51^\circ\,31'\,45''{,}21$ mit den beiden Abweichungen von $4''{,}1$ auf S.~\href{https://api.digitale-sammlungen.de/iiif/image/v2/bsb10538349_00208/full/1600,/0/default.jpg}{204} gehören zu diesem Briefpaar. Der eigenständige Julibrief hat seine eigenen Lesarten. KI-Übertragung und KI-Satz.
\clearpage
